\chapter{KẾT LUẬN}
\label{ch:ket-luan}

\section{Tóm tắt nội dung thực hiện}

Đồ án đã xác định quy trình dự báo xu hướng biến động giá cổ phiếu trên sàn HOSE. Trong
quy trình này, tin tức tài chính tiếng Việt được xử lý bằng PhoBERT, còn dữ liệu giá được
biểu diễn bằng lợi suất và các đặc trưng kỹ thuật. Hai nhánh dữ liệu được căn chỉnh theo
mã cổ phiếu, ngày giao dịch và mốc cắt thông tin, sau đó được đưa vào thang mô hình từ
đối chứng cơ sở đến LSTM hai nhánh có cơ chế hợp nhất. Bộ biến đổi hợp nhất theo thời gian được
đặt ở vị trí mở rộng để không làm thay đổi cấu hình cơ sở của đồ án.

Quy trình gồm thu thập dữ liệu, kiểm tra thông tin truy nguyên và bản ghi trùng, gán nhãn
cảm xúc, sinh xác suất của ba lớp, tổng hợp đặc trưng theo mã và ngày, tạo nhãn UP, FLAT
và DOWN, huấn luyện, đánh giá cuốn chiếu theo thời gian và phân tích loại bỏ đặc trưng.
Tiền xử lý của nhánh giá chỉ học từ tập huấn luyện, mốc cắt thông tin là 15 giờ và năm
đoạn kiểm thử cuốn chiếu không chồng lấn. Thang dự báo chính dùng điểm kiểm cảm xúc học từ
nhãn gồm cả giai đoạn kiểm thử nên là lượt hồi cứu; một lượt bổ sung với điểm kiểm đóng
băng trước ngày kiểm thử đầu tiên cung cấp đánh giá ngoài mẫu hợp lệ.

\section{Đối chiếu với mục tiêu và câu hỏi nghiên cứu}

Mục tiêu thứ nhất là tạo bộ dữ liệu có thể truy nguyên. Mục tiêu này \textbf{đã hoàn
thành}: 14.256 liên kết tin--mã Vietstock với 100\% có mốc thời gian và nội dung, và
14.990 dòng giá của 10 mã HOSE trong giai đoạn 2020--2025 không có giá trị thiếu. Trong
12.624 liên kết thuộc cửa sổ nghiên cứu, 12.613 được neo vào một phiên giao dịch và 11 bị
loại vì không neo được. Toàn bộ dữ liệu được ghi kèm mã băm trong bản kê của lần chạy.

Mục tiêu thứ hai là đánh giá PhoBERT ở tầng cảm xúc. Xác thực chéo theo tầng trên 1.306
nhãn đã rà soát đã hoàn thành: cấu hình có trọng số nghịch đảo tần suất đạt F1 vĩ mô
0,7298, độ chính xác 0,8125 và độ chính xác cân bằng 0,7689. Con số 0,8125 vẫn dưới mốc
tham chiếu 85--90\% của đề cương. Năm tập kiểm thử ngoài có tổng 59 mẫu NEGATIVE và độ
bao phủ trung bình 0,6803, nên lớp này đã được đo nhưng bất định vẫn lớn hơn hai lớp
còn lại. Kết quả có trọng số là phép so sánh trên tập mở rộng, không phải mô hình cuối.
Sau khi khóa cấu hình, checkpoint triển khai đã được tinh chỉnh lại trên toàn bộ 1.306
nhãn với năm epoch cố định, chấm lại 12.624 liên kết tin--mã và dùng làm đầu vào của dự báo.
Lượt tinh chỉnh lại không tạo chỉ số kiểm thử mới.

Mục tiêu thứ ba và thứ tư là kết hợp hai nguồn dữ liệu và huấn luyện mô hình dự báo. Mục
tiêu này \textbf{đã hoàn thành với cả LSTM và TFT}. Mô hình LSTM chỉ giá đạt F1 vĩ mô
0,3839, độ chính xác cân bằng 0,3918 và AUC-ROC một-chống-còn-lại 0,5728; các điểm này cao
hơn mốc ngẫu nhiên 0,3346 và mốc lớp phổ biến 0,1573. TFT đạt 0,3595 khi chỉ dùng giá,
thấp hơn LSTM, nên kiến trúc phức tạp hơn không tạo lợi thế trên quy mô dữ liệu này. Báo
cáo không có kiểm định ghép cặp với hai mốc đối chứng, nên so sánh này chỉ mô tả mức điểm
quan sát được \cite{Halder2022_221107392v1,Gu2024_240716150v1}.

Mục tiêu thứ năm là đo phần đóng góp riêng của tín hiệu cảm xúc. Phép đo ghép cặp \textbf{đã
hoàn thành ở dạng hồi cứu}. Phép đo được thực hiện trên cùng cửa sổ, cùng ba hạt giống,
cùng quy trình tiền xử lý và cùng tập dòng kiểm thử. Đối chứng giữ nguyên kiến trúc hai
nhánh, cùng bài tin, thời điểm và số lượng tin, nhưng thay xác suất cảm xúc bằng phân phối
tiên nghiệm trung tính. Kết quả:

\begin{itemize}[nosep]
  \item hiệu số trực tiếp: $+0{,}0110$ F1 vĩ mô;
  \item trong đó hiệu ứng kiến trúc, sự hiện diện và khối lượng tin: $+0{,}0044$;
  \item đóng góp thông tin phân cực: $+0{,}0066$ theo cửa sổ; ước lượng theo ngày là
        $+0{,}0014$, khoảng tin cậy 95\% $[-0{,}0087; +0{,}0116]$.
\end{itemize}

Khoảng tin cậy theo ngày chứa 0 ở cả ba chỉ số, và bộ sinh cảm xúc của lượt này đã học từ
nhãn trong giai đoạn kiểm thử, nên các số liệu trên chỉ là chênh lệch hồi cứu.

Để có đánh giá ngoài mẫu hợp lệ, một điểm kiểm thứ hai được tinh chỉnh chỉ trên 1.049 nhãn
công bố trước ngày quan sát kiểm thử đầu tiên (22/10/2024), rồi chấm lại toàn bộ 14.256
liên kết tin--mã và chạy lại thang dự báo cùng đối chứng ghép cặp. Với điểm kiểm đóng băng,
đóng góp thông tin phân cực là $-0{,}0024$ F1 vĩ mô ở LSTM hai nhánh (khoảng tin cậy theo
ngày $[-0{,}0134; +0{,}0079]$) và $+0{,}0061$ ở TFT ($[-0{,}0041; +0{,}0181]$); cả hai
khoảng đều chứa 0. So với lượt hồi cứu, ước lượng điểm giảm và đổi dấu ở LSTM, phù hợp với
giả thuyết rằng một phần hiệu số dương trước đó đến từ rò rỉ thời điểm qua bộ sinh.

Kết luận ngoài mẫu của đồ án vì vậy là: trên bộ dữ liệu và thiết lập hiện tại, cảm xúc tin
tài chính tiếng Việt \textbf{không mang lại cải thiện có thể tái lập} cho bài toán phân loại
xu hướng phiên kế tiếp. Đây là kết quả âm tính được đo đúng quy trình, không phải thất bại
của quy trình đo.

Mục tiêu thứ sáu là hệ thống demo hoàn chỉnh. Mục tiêu này \textbf{đã hoàn thành một phần}:
đường dự báo hằng ngày (thu giá, thu tin, chấm cảm xúc, dự báo, đối soát kết quả thực tế)
chạy qua dòng lệnh và có bản dashboard đọc artifact, phục vụ tín hiệu mới nhất của năm cấu
hình mô hình cùng lịch sử độ chính xác đã đối soát. Giao diện trực quan đầy đủ cho người
dùng cuối vẫn là hướng phát triển.

Câu hỏi thứ ba của đề cương --- điều kiện nào cảm xúc phát huy tác dụng --- chưa được trả
lời bằng ước lượng ghép cặp phân tầng. Đối chứng mới đã kiểm soát sự hiện diện và khối
lượng tin trên toàn bộ mẫu; phân tầng theo giai đoạn thị trường, ngày sự kiện và mật độ tin
chưa được thực hiện.

\section{Đóng góp và ý nghĩa}

Đóng góp chính của đồ án là một quy trình thực nghiệm đa phương thức cho dữ liệu tiếng
Việt được mô tả tường minh. Quy trình nối mô hình ngôn ngữ mở với nhánh chuỗi giá, đồng
thời đặt phép loại bỏ đặc trưng và kiểm soát tại thời điểm quan sát ở vị trí trung tâm.
Điều này cần thiết vì các hướng hợp nhất dùng biểu diễn ngữ cảnh, sự kiện hoặc mô hình
ngôn ngữ lớn có năng lực biểu đạt cao hơn nhưng cũng đòi hỏi nhiều dữ liệu, chi phí và
điều kiện tái lập hơn \cite{Chen2021_210708721v1,Elahi2024_241101368v1,Liang2024_241210823v2}.

Về phương pháp, đồ án phân biệt ba loại kết luận. Loại thứ nhất là chất lượng phân loại
cảm xúc trên tập được gán nhãn. Loại thứ hai là khả năng dự báo ngoài mẫu của từng mô
hình. Loại thứ ba là phần đóng góp biên của tín hiệu cảm xúc qua phép loại bỏ đặc trưng.
Ba loại kết luận này có liên quan nhưng không thể thay thế cho nhau.

Về thực tiễn, quy trình có thể làm nền cho công cụ thăm dò dữ liệu, giúp người nghiên cứu
quan sát phân bố tin, xác suất cảm xúc, trạng thái kỹ thuật và lớp dự báo theo ngày. Tuy
nhiên, quy trình chưa phải là hệ thống giao dịch và không cung cấp khuyến nghị đầu tư.
Giá trị kinh tế sau phí giao dịch, trượt giá và giới hạn thanh khoản cần được đánh giá
trong một nghiên cứu riêng.

\section{Giới hạn của đồ án}

Thứ nhất, phạm vi chỉ gồm 10 mã VN30 trên HOSE trong giai đoạn 2020 đến hết năm 2025.
Đặc điểm ngành, thanh khoản và mật độ tin của các mã này có thể ảnh hưởng đến kết quả,
nên chưa đủ cơ sở để suy rộng sang cổ phiếu vốn hóa nhỏ, HNX, UPCOM hoặc thị trường
quốc tế.

Thứ hai, dữ liệu tin phụ thuộc vào nguồn, chính sách truy cập và chất lượng thời điểm.
Bài viết bị thiếu thời gian hoặc được cập nhật nhiều lần có thể làm thay đổi cách ánh xạ.
Mốc cắt thông tin bảo thủ giúp giảm nguy cơ sử dụng dữ liệu tương lai, nhưng cũng có thể
làm trễ một phần tín hiệu được thị trường nhìn thấy trước phiên kế tiếp.

Thứ ba, và là giới hạn quyết định: \textbf{điểm kiểm cảm xúc được tinh chỉnh trên toàn bộ
1.306 nhãn, gồm cả bài sau các ngày kiểm thử dự báo}. Các xác suất dùng trong nhánh có tin
vì vậy có thể chịu ảnh hưởng gián tiếp từ dữ liệu tương lai. Cần tinh chỉnh lại trong từng
cửa sổ chỉ bằng nhãn quá khứ, hoặc dùng một điểm kiểm đóng băng trước ngày kiểm thử đầu
tiên. Chỉ sau đó mới có thể dùng khoảng tin cậy để kết luận về hiệu quả ngoài mẫu. Ngoài
ra, nhãn ba lớp buộc phải nén các bài có nhiều sự kiện trái chiều về một lớp duy nhất.

Thứ tư, dự báo phiên kế tiếp chỉ quan sát một khoảng thời gian ngắn. Tin về kế hoạch dài
hạn, chất lượng quản trị hoặc thay đổi ngành có thể chưa thể hiện trong lợi suất của phiên
sau. Ngoài ra, chỉ số phân loại không phản ánh trực tiếp khả năng sinh lợi. Các nghiên cứu
dùng mô phỏng giao dịch hoặc mô hình quan hệ liên mã cho thấy cần thêm nhiều góc nhìn,
nhưng những mở rộng đó nằm ngoài cấu hình cơ sở
\cite{Chen2021_210708721v1,Xu2021_211013716v2}.

Thứ năm, ngưỡng phân vị một phần ba định nghĩa lớp FLAT theo nghĩa thống kê chứ không theo
nghĩa kinh tế. Ngưỡng khoảng $\pm0{,}5\%$ là một lựa chọn thiết kế; biến thể dải cố định
chưa được kiểm tra. Cũng vì các lớp gần cân bằng theo thiết kế, mức ngẫu nhiên của bài toán
là 0,333 chứ không phải 0,5 như cách phát biểu hai lớp trong đề cương.

Thứ sáu, cả hai mô hình LSTM nghiêng về lớp FLAT: độ bao phủ FLAT khoảng 0,61--0,65 trong khi
độ bao phủ DOWN chỉ khoảng 0,27. Phần lớn giá trị F1 vĩ mô đến từ lớp giữa, là lớp ít hữu ích
nhất cho người sử dụng.

\section{Kết luận cuối cùng}

Đề tài đưa ra một cách tiếp cận có kiểm soát để kiểm tra việc bổ sung cảm xúc tin tức
tiếng Việt vào dự báo xu hướng giá cổ phiếu. PhoBERT được lựa chọn làm mô hình ngôn ngữ,
LSTM hai nhánh được lựa chọn làm cấu hình cơ sở để so sánh mô hình chỉ dùng giá với mô
hình có cảm xúc, và TFT được chạy như tầng mở rộng sau khi cấu hình cơ sở được xác minh.

Đồ án \textbf{đã đưa ra kết quả định lượng ở cả hai chế độ}. Mô hình dự báo dùng đặc trưng
giá có điểm cao hơn hai mốc đối chứng bắt buộc; do chưa có kiểm định ghép cặp với các mốc
này, kết quả chỉ mô tả mức điểm quan sát được. Với đối chứng giữ nguyên kiến trúc, thời
điểm và số lượng tin, lượt hồi cứu cho đóng góp thông tin theo cửa sổ $+0{,}0066$ F1 vĩ mô
(theo ngày $+0{,}0014$, khoảng tin cậy $[-0{,}0087; +0{,}0116]$).

Lượt ngoài mẫu với điểm kiểm đóng băng trước ngày kiểm thử đầu tiên cho $-0{,}0024$ F1 vĩ
mô ở LSTM hai nhánh (khoảng tin cậy theo ngày $[-0{,}0134; +0{,}0079]$) và $+0{,}0061$ ở
TFT ($[-0{,}0041; +0{,}0181]$). Mọi khoảng tin cậy đều chứa 0, và ước lượng điểm của LSTM
đổi dấu khi loại bỏ rò rỉ thời điểm.

Kết luận của đồ án vì vậy là: trên 10 mã HOSE giai đoạn 2020--2025 với nguồn tin Vietstock
và chân trời một phiên, cảm xúc tin tài chính tiếng Việt \textbf{không mang lại cải thiện
có thể tái lập} cho bài toán phân loại xu hướng. Giá trị của đồ án nằm ở chỗ kết luận âm
tính này được rút ra từ một quy trình đo có đối chứng ghép cặp giữ nguyên kiến trúc và độ
phủ tin, có điểm kiểm phù hợp thời điểm, và có đủ bản kê để tái lập --- thay vì một con số
dương thu được từ thiết lập bị rò rỉ.

Mọi diễn giải phải giữ đúng phạm vi 10 mã HOSE, giai đoạn 2020--2025, nguồn tin Vietstock và
chân trời một phiên. Kết luận không được suy rộng thành lời hứa hiệu năng ngoài dữ liệu, và
không phải khuyến nghị đầu tư.
